\documentclass[leqno]{article}
% Shared mark scheme preamble for resources.danieljsmith.org: the pack
% preamble, unchanged, then the scheme's own layer. A scheme's .tex is
% rendered from its YAML (scripts/README.md, Mark schemes), never edited.
\input{\CurrentFilePath/pack_preamble.tex}

% The page: equal side margins, so the table runs the full width of the
% header rule. The footers' right edge follows the text's.
\geometry{left=0.8in, right=0.8in}
\makeatletter
\renewcommand{\djs@footR}[1]{%
  \put(\dimexpr1in+\hoffset+\oddsidemargin+\textwidth\relax,
    -\dimexpr\paperheight-0.32in\relax){\makebox[0pt][r]{#1}}}
\makeatother

\usepackage[export]{adjustbox}
\usepackage{tabularray}
\definecolor{msGrid}{RGB}{176, 172, 164}
\definecolor{msTint}{RGB}{245, 243, 239}

% The header on every page, the front page included, and the booklet label
% under the contents-page title. \djsQuestionPageStyle stays empty.
\newcommand{\djsSchemeHeader}[1]{%
  \djsPackHeader[markscheme]{\textit{Mark scheme} - #1}%
  \renewcommand{\djsPackTitle}{#1}%
  \renewcommand{\djsBookletLabel}{Mark scheme}%
}

% The front page: the contents heading, spaced as tocloft spaces the packs'
% "Questions", then the contents table. A question's name and page number
% both link to its first page.
\newcommand{\msContentsHeading}{%
  \vskip\cftbeforetoctitleskip
  {\noindent\cfttoctitlefont\contentsname\par}%
  \vskip\cftaftertoctitleskip}
\NewTblrEnviron{mscontents}
\SetTblrInner[mscontents]{
  colspec = {Q[wd=26mm, l] X[l] Q[wd=16mm, c] Q[wd=14mm, r]},
  width = \linewidth,
  rowsep = 4pt, colsep = 6pt,
  hlines = {0.4pt, msGrid}, vlines = {0.4pt, msGrid},
  hline{2} = {2.2pt, scarlet},
  row{1} = {font=\small\bfseries},
  row{Z} = {font=\bfseries},
  column{2} = {font=\small},
}
\newcommand{\msContentsQuestion}[1]{\hyperref[q:#1]{Question #1}}
\newcommand{\msContentsPage}[1]{\hyperref[q:#1]{\pageref*{q:#1}}}
\newcommand{\msSep}{\hspace{0.9em}}

% The conventions pages' layout: headings in the documents' own bold, tables
% ruled and padded as the scheme's, roomy lists.
\newcommand{\msHead}[1]{%
  \par\vspace{18pt}%
  {\noindent\bfseries #1\par}%
  \vspace{6pt}}
\NewTblrEnviron{mstblr}
\SetTblrInner[mstblr]{rowsep = 6pt, colsep = 6pt,
  hlines = {0.4pt, msGrid}, vlines = {0.4pt, msGrid}}
\newenvironment{msList}
  {\begin{itemize}[leftmargin=1.2em, itemsep=5pt, topsep=3pt, parsep=0pt]}
  {\end{itemize}}

% A question's first page: the link target and the PDF bookmark, as
% \questionitem gives the packs.
\newcommand{\msQnum}{}
\newcommand{\msQuestion}[1]{%
  \clearpage
  \renewcommand{\msQnum}{#1}%
  \phantomsection\label{q:#1}%
  \addcontentsline{toc}{subsection}{\protect\djsTocEntry{#1}}}

% The scheme's table: Question 16mm, Scheme the rest, Marks 16mm, Notes
% 44mm; every cell ruled in msGrid; the scarlet rule under the head row. A
% long table has no caption or "continued" foot; a continuation page opens
% with a muted line above the repeated head row. The rules across a part's
% rows, the alternative's tint and the total row are the rendered table's.
\NewTblrTheme{ms}{%
  \DefTblrTemplate{firsthead}{default}{}%
  \DefTblrTemplate{middlehead,lasthead}{default}{%
    {\small\itshape\color{djsMuted}Question \msQnum\ continued\par}%
    \vspace{6pt}}%
  \DefTblrTemplate{firstfoot,middlefoot,lastfoot}{default}{}%
}
\NewTblrEnviron{msscheme}
\SetTblrOuter[msscheme]{long, theme=ms, entry=none, label=none}
\SetTblrInner[msscheme]{
  colspec = {Q[wd=16mm, l] X[l] Q[wd=16mm, c] Q[wd=44mm, l]},
  width = \linewidth,
  rowsep = 6pt, colsep = 6pt,
  rowhead = 1,
  vlines = {0.4pt, msGrid},
  hline{1} = {0.4pt, msGrid}, hline{2} = {2.2pt, scarlet},
  hline{Z} = {0.4pt, msGrid},
  column{1} = {font=\bfseries}, column{3} = {font=\bfseries},
  column{4} = {font=\small},
  row{1} = {rowsep=4pt, font=\small\bfseries},
}

% Two marks on a row: each note starts with its code.
\newcommand{\msNote}[2]{\textbf{#1}\enspace #2}
% A full-row figure: the worked solution's own picture, scaled uniformly and
% only down, to the table less the Question column. The picture is drawn
% with the line width it has in the solutions document, at its list depth
% there (the optional argument, 1 by default: a part's working), so one
% sized by \linewidth keeps its size; tabularray sets the cell's \linewidth
% to the cell's measured width, which would shrink it at every pass. The
% widths are measured once, in a solution box at the packs' text width
% (pack_preamble.tex's margins, 0.8in and 1in).
\newlength{\msSpanW}
\makeatletter
\newcommand{\ms@keepLW}[1]{%
  \expandafter\xdef\csname ms@solnLW@#1\endcsname{\the\linewidth}}
\AtBeginDocument{%
  \setlength{\msSpanW}{\dimexpr\textwidth-16mm-28pt\relax}%
  \setbox\z@\vbox{%
    \begin{minipage}{\dimexpr\paperwidth-1.8in\relax}
      \begin{djsSolution}
        \ms@keepLW{0}%
        \begin{enumerate}\item \ms@keepLW{1}%
        \begin{enumerate}\item \ms@keepLW{2}%
        \begin{enumerate}\item \ms@keepLW{3}%
        \end{enumerate}\end{enumerate}\end{enumerate}
      \end{djsSolution}
    \end{minipage}}}
\newcommand{\msFigure}[2][1]{%
  \adjustbox{max width=\msSpanW, max totalheight=0.6\textheight, valign=t}{%
    \hsize=\csname ms@solnLW@#1\endcsname
    \linewidth=\hsize
    #2}}
\makeatother
% A figure mark's feature, in the Scheme column.
\newcommand{\msFeature}[1]{\textit{#1}}
% An alternative's title row.
\newcommand{\msAlternative}[1]{\textit{Alternative:} #1}
% Drafts only: a queried mark's code, the query's why under its note, and a
% flag in its part's place.
\newcommand{\msQueried}[2]{#1~\(\to\) #2?}
\newcommand{\msQuery}[1]{{\itshape\color{djsMuted}Query: #1}}
\newcommand{\msFlag}[2]{\textbf{#1:}\enspace #2}


\djsSchemeHeader{Roots of Polynomials}
\djsSitePath{Core Pure}{Miscellaneous Pure}

\djsLastUpdated{24 October 2026}
\begin{document}
\djsSetQuestionCount{15}
\djsAddMarks{1}{5}
\djsAddMarks{2}{8}
\djsAddMarks{3}{7}
\djsAddMarks{4}{6}
\djsAddMarks{5}{6}
\djsAddMarks{6}{8}
\djsAddMarks{7}{4}
\djsAddMarks{8}{10}
\djsAddMarks{9}{7}
\djsAddMarks{10}{11}
\djsAddMarks{11}{10}
\djsAddMarks{12}{7}
\djsAddMarks{13}{10}
\djsAddMarks{14}{10}
\djsAddMarks{15}{6}
\vspace*{6pt}
\djsContentsTitle
\vspace{30pt}
\msContentsHeading
\begin{mscontents}{}
Question & Marks by part & Marks & Page \\
\msContentsQuestion{1} & (a)\,1\msSep (b)\,1\msSep (c)\,3 & 5 & \msContentsPage{1} \\
\msContentsQuestion{2} & (a)\,1\msSep (b)(i)\,2\msSep (b)(ii)\,2\msSep (b)(iii)\,3 & 8 & \msContentsPage{2} \\
\msContentsQuestion{3} & (a)\,4\msSep (b)\,3 & 7 & \msContentsPage{3} \\
\msContentsQuestion{4} & & 6 & \msContentsPage{4} \\
\msContentsQuestion{5} & (a)\,1\msSep (b)\,2\msSep (c)\,3 & 6 & \msContentsPage{5} \\
\msContentsQuestion{6} & (a)\,3\msSep (b)\,3\msSep (c)\,2 & 8 & \msContentsPage{6} \\
\msContentsQuestion{7} & & 4 & \msContentsPage{7} \\
\msContentsQuestion{8} & (a)\,4\msSep (b)\,3\msSep (c)\,3 & 10 & \msContentsPage{8} \\
\msContentsQuestion{9} & (a)\,6\msSep (b)\,1 & 7 & \msContentsPage{9} \\
\msContentsQuestion{10} & (a)\,3\msSep (b)\,3\msSep (c)\,5 & 11 & \msContentsPage{10} \\
\msContentsQuestion{11} & (a)\,4\msSep (b)\,3\msSep (c)\,3 & 10 & \msContentsPage{11} \\
\msContentsQuestion{12} & (a)\,3\msSep (b)\,4 & 7 & \msContentsPage{12} \\
\msContentsQuestion{13} & (a)\,6\msSep (b)\,4 & 10 & \msContentsPage{13} \\
\msContentsQuestion{14} & (a)\,5\msSep (b)\,3\msSep (c)\,2 & 10 & \msContentsPage{14} \\
\msContentsQuestion{15} & (a)\,3\msSep (b)\,3 & 6 & \msContentsPage{15} \\
\SetCell[c=2]{r} Total & & 115 & \\
\end{mscontents}
\newpage
% The conventions pages, shared by every pack's mark scheme: one source of
% truth, \input after the front page, with no strand or pack logic. The keys
% on the first side, the rules on the second. The scheme layer defines
% \msHead, msList and mstblr.
\vspace*{6pt}%
{\noindent\cfttoctitlefont Reading the mark scheme\par}%
\vspace{10pt}%
These rules apply throughout, unless a note says otherwise.

\msHead{Marks}
Every mark is worth one, so \textbf{M1 A1} on a row is two marks. Only marks
coded ft follow through.\par
\vspace{10pt}%
\begin{mstblr}{
  colspec = {Q[wd=15mm, l] X[l]},
  width = \linewidth,
  hline{2} = {2.2pt, scarlet},
  column{1} = {font=\bfseries},
  row{1} = {font=\small\bfseries, rowsep=4pt},
}
Code & Meaning \\
M1 & Method: a correct method, applied to this question. A slip doesn't cost it if the method is clear. A formula must be used, not just quoted. \\
dM1 & Dependent method: as M1, but only if the method mark before it was earned, unless its note names another. \\
A1 & Accuracy: a correct answer or step. It needs the method mark it follows. \\
A1ft & As A1, following through from an earlier answer that the note names. \\
A1~AG & The last mark of a part whose answer is printed in the question: complete working, with no errors. \\
B1 & A correct value, fact or statement. It needs no method mark. \\
B1ft & As B1, following through from an earlier answer that the note names. \\
E1 & A written explanation, interpretation or conclusion. Concluding a proof by induction needs every earlier mark. \\
\end{mstblr}

\msHead{Abbreviations}
ft and AG qualify a code in the Marks column. The others appear in the
notes.\par
\vspace{10pt}%
\begin{mstblr}{
  colspec = {Q[wd=17mm, l] X[l] Q[wd=15mm, l] X[l]},
  width = \linewidth,
  column{1} = {font=\bfseries}, column{3} = {font=\bfseries},
}
ft & follow through & AG & answer given in the question \\
cao & correct answer only, with no follow-through & cso & correct solution only: no errors in the part \\
oe & or an equivalent form & awrt & any answer that rounds to this value \\
soi & seen, or implied by later working & isw & ignore working after a correct answer \\
SC & special case: credit outside the main method & their & the student's own earlier value \\
condone & a slip that doesn't cost the mark & [\;\,] & working the mark doesn't need \\
s.f. & significant figures & d.p. & decimal places \\
\end{mstblr}

\newpage
\msHead{Reading the scheme}
\begin{msList}
\item Each row is a line of working, with the marks it earns beside it.
  Where a row carries two marks, each note starts with its code.
\item \textbf{[\textit{n}]} ends each part with its total.
\item A shaded \textit{Alternative} is a second method for the same total.
  Mark an answer by the method it follows.
\item Any other valid method earns the equivalent marks, unless the question
  names a method or says \textit{hence}.
\item A sketch earns one B1 for each feature named below its figure.
\end{msList}

\msHead{Accuracy and answers}
\begin{msList}
\item Unless a question says otherwise, give answers to 3 significant
  figures, and angles in degrees to 1 decimal place. Angles in radians take
  3 significant figures.
\item A final answer with more figures is accepted if it rounds to the
  scheme's value; one with fewer is not. An answer spoiled by rounding too
  early loses its accuracy mark.
\item Where a question asks for an exact answer, a decimal loses the final
  accuracy mark. Otherwise either is accepted.
\item Wrong or missing units cost nothing unless the question asks for them.
\item A correct answer with no working earns full marks, unless the question
  needs the working: show that, prove, hence, a method it names, or an exact
  answer. A wrong answer with no working earns nothing. Calculator answers
  are treated the same way.
\item A correct answer, once reached, isn't spoiled by later working (isw).
\end{msList}

\msHead{Value of $\boldsymbol{g}$}
\begin{msList}
\item Take $g = 9.8\,\mathrm{m\,s^{-2}}$ unless a question states a value,
  which is then the only one accepted. Answers that use $g = 9.8$ may be
  given to 2 or 3 significant figures.
\item Where a question doesn't state $g$, an answer worked with $g = 10$
  throughout earns the same marks, keeping any value printed in the
  question. The worked solutions list the final answers with $g = 10$.
\end{msList}


\msQuestion{1}
\begin{msscheme}{
  hline{3,5,7-8} = {2-Z}{0.4pt, msGrid},
  hline{4,6,9} = {0.4pt, msGrid},
  row{9} = {font=\bfseries},
}
Question & Scheme & Marks & Notes \\
1(a) & {\(p = -6\)} & B1 & \\
\nopagebreak & & [1] & \\
1(b) & {\(\alpha\beta+\beta\gamma+\gamma\alpha = -4\), \(\alpha\beta\gamma = \frac83\)} & B1 & {Both values needed.} \\
\nopagebreak & & [1] & \\
1(c) & {\(\frac{(\alpha+1)(\beta+1)(\gamma+1)}{\alpha\beta\gamma}\)\\[4pt] \(= \frac{\alpha\beta\gamma + (\alpha\beta+\beta\gamma+\gamma\alpha) + (\alpha+\beta+\gamma) + 1}{\alpha\beta\gamma}\)} & M1 A1 & {\msNote{M1}{Writes the product over \(\alpha\beta\gamma\), or expands it directly, giving terms in the root sums. Condone one slip in the expansion.}\\ \msNote{A1}{Correct expression in \(\alpha+\beta+\gamma\), \(\alpha\beta+\beta\gamma+\gamma\alpha\) and \(\alpha\beta\gamma\), oe such as \(1+\sum\frac1\alpha+\sum\frac1{\alpha\beta}+\frac1{\alpha\beta\gamma}\) rewritten in them.}} \\
& {\(= \frac{\frac83 - 4 + 2 + 1}{\frac83} = \frac58\)} & A1 & {\(\frac58\) oe, such as 0.625.} \\
\nopagebreak & & [3] & \\
\nopagebreak \SetCell[c=2]{r} Total for Question 1 & & 5 & \\
\end{msscheme}

\msQuestion{2}
\begin{msscheme}{
  hline{3,5-6,8-9,11-13} = {2-Z}{0.4pt, msGrid},
  hline{4,7,10,14} = {0.4pt, msGrid},
  row{14} = {font=\bfseries},
}
Question & Scheme & Marks & Notes \\
2(a) & {\(\alpha+\beta+\gamma=-5,\ \alpha\beta+\beta\gamma+\gamma\alpha=-4,\ \alpha\beta\gamma=24\)} & B1 & \\
\nopagebreak & & [1] & \\
2(b)(i) & {\(\frac{\alpha+\beta+\gamma}{\alpha\beta\gamma}=\frac{\their{-5}}{\their{24}}\)} & M1 & {Combines over the common denominator \(\alpha\beta\gamma\) to the sum of roots over the product and substitutes their values from (a).} \\
& {\(-\frac{5}{24}\)} & A1 & {oe, or awrt \(-0.208\).} \\
\nopagebreak & & [2] & \\
2(b)(ii) & {\(\alpha\beta\gamma+2(\alpha\beta+\beta\gamma+\gamma\alpha)+4(\alpha+\beta+\gamma)+8\)\\[4pt] \(\phantom{=} = \their{24}+2(\their{-4})+4(\their{-5})+8\)} & M1 & {Expands to the three root sums and a constant, or uses \(x=y-2\), and substitutes their values. Condone one coefficient slip.} \\
& {\(4\)} & A1 & \\
\nopagebreak & & [2] & \\
2(b)(iii) & {\(\alpha^2\beta^2+\beta^2\gamma^2+\gamma^2\alpha^2\)\\[4pt] \(\phantom{=} = (\alpha\beta+\beta\gamma+\gamma\alpha)^2-2\alpha\beta\gamma(\alpha+\beta+\gamma)\)} & M1 & {Writes the sum as \((\alpha\beta+\beta\gamma+\gamma\alpha)^2\) minus a multiple of \(\alpha\beta\gamma(\alpha+\beta+\gamma)\). Condone a wrong coefficient.} \\
& {\(\phantom{=} = (-4)^2-2(24)(-5)\)} & A1 & {Correct identity with the correct values substituted.} \\
& {\(\phantom{=} = 256\)} & A1 & \\
\nopagebreak & & [3] & \\
\nopagebreak \SetCell[c=2]{r} Total for Question 2 & & 8 & \\
\end{msscheme}

\msQuestion{3}
\begin{msscheme}{
  hline{3-6,8-10} = {2-Z}{0.4pt, msGrid},
  hline{7,11} = {0.4pt, msGrid},
  row{11} = {font=\bfseries},
}
Question & Scheme & Marks & Notes \\
3(a) & {\(\alpha+\beta+\gamma=-2,\ \alpha\beta+\beta\gamma+\gamma\alpha=-5,\ \alpha\beta\gamma=3\)} & B1 & {All three, soi.} \\
& {\((\alpha\beta+\beta\gamma+\gamma\alpha)^2\)\\[4pt] \(\phantom{(} = \sum\alpha^2\beta^2 + 2(\alpha\beta^2\gamma+\beta\gamma^2\alpha+\gamma\alpha^2\beta)\)} & M1 & {Squares \(\sum\alpha\beta\), giving the three squared terms and cross terms. Condone a slip in the cross terms' coefficient.} \\
& {\(\sum\alpha^2\beta^2 = \left(\sum\alpha\beta\right)^2 - 2\alpha\beta\gamma\sum\alpha\)} & A1 & {Correct identity in terms of \(\sum\alpha\beta\), \(\alpha\beta\gamma\) and \(\sum\alpha\), oe.} \\
& {\(= (-5)^2 - 2(3)(-2) = 37\)} & A1 & {cso: withhold if sign errors in the root sums cancel.} \\
\nopagebreak & & [4] & \\
3(b) & {\(\sum\alpha^2 = (-2)^2 - 2(-5)\ [= 14]\)} & M1 & {Uses \(\left(\sum\alpha\right)^2 - 2\sum\alpha\beta\) with their values. Condone sign slips in substitution.} \\
& {\(\alpha^2\beta^2\gamma^2 = 3^2 = 9\)} & B1 & \\
& {\(x^3 - 14x^2 + 37x - 9 = 0\)} & A1 & {Or any integer multiple, in any variable; must be an equation.} \\
\nopagebreak & & [3] & \\
\nopagebreak \SetCell[c=2]{r} Total for Question 3 & & 7 & \\
\end{msscheme}

\msQuestion{4}
\begin{msscheme}{
  hline{3-6} = {2-Z}{0.4pt, msGrid},
  hline{7} = {0.4pt, msGrid},
  row{7} = {font=\bfseries},
}
Question & Scheme & Marks & Notes \\
4 & {\((a-d)+a+(a+d)=-3 \Rightarrow a=\ldots\)} & M1 & {Equates the sum of equally spaced roots, such as \(a-d,\ a,\ a+d\) or \(a,\ a+d,\ a+2d\), to \(\pm3\).} \\
& {\(a=-1\)} & A1 & {Middle root \(-1\), in any parametrisation.} \\
& {\((a-d)a(a+d)=20\)} & M1 A1 & {\msNote{M1}{Equates the product of their three roots to \(\pm20\).}\\ \msNote{A1}{Correct equation, oe such as \(a(a^2-d^2)=20\) or with \(a=-1\) substituted.}} \\
& {\(-1\left(1-d^2\right)=20 \Rightarrow d^2=\ldots\ [=21]\)} & dM1 & {Substitutes their \(a\) and solves for \(d\) or \(d^2\).} \\
& {\(x=-1-\sqrt{21},\ -1,\ -1+\sqrt{21}\)} & A1 & {All three, exact or awrt \(-5.58\) and \(3.58\), with no others.} \\
\nopagebreak \SetCell[c=2]{r} Total for Question 4 & & 6 & \\
\end{msscheme}

\msQuestion{5}
\begin{msscheme}{
  hline{3,5-6,8-10} = {2-Z}{0.4pt, msGrid},
  hline{4,7,11} = {0.4pt, msGrid},
  row{11} = {font=\bfseries},
}
Question & Scheme & Marks & Notes \\
5(a) & {\(\alpha + \beta = k,\quad \alpha\beta = 1\)} & B1 & {Both values needed.} \\
\nopagebreak & & [1] & \\
5(b) & {\(\alpha^2 + \beta^2 + \alpha + \beta\)\\[4pt] \(\phantom{\alpha^2} = (\alpha+\beta)^2 - 2\alpha\beta + (\alpha+\beta)\)} & M1 & {Writes \(\alpha^2+\beta^2\) in terms of \(\alpha+\beta\) and \(\alpha\beta\) and substitutes their values from (a). Condone a sign slip in the identity.} \\
& {\(= k^2 + k - 2\)} & A1 & {Must be simplified.} \\
\nopagebreak & & [2] & \\
5(c) & {\((\alpha^2+\beta)(\beta^2+\alpha)\)\\[4pt] \(\phantom{(\alpha} = (\alpha\beta)^2 + \alpha^3 + \beta^3 + \alpha\beta\)\\[4pt] \(\alpha^3+\beta^3 = (\alpha+\beta)^3 - 3\alpha\beta(\alpha+\beta)\)} & M1 & {Expands the product to four terms and expresses \(\alpha^3+\beta^3\) in terms of \(k\) by any valid identity or reduction. Condone sign and coefficient slips.} \\
& {Product \(= k^3 - 3k + 2\)} & A1 & \\
& {\(x^2 - (k^2+k-2)x + k^3 - 3k + 2 = 0\)} & A1 & {Or any integer multiple, in \(x\), with \(= 0\); brackets need not be expanded. Stating \(p\), \(q\), \(r\) alone scores A0.} \\
\nopagebreak & & [3] & \\
\nopagebreak \SetCell[c=2]{r} Total for Question 5 & & 6 & \\
\end{msscheme}

\msQuestion{6}
\begin{msscheme}{
  hline{3-5,7-9,11-12} = {2-Z}{0.4pt, msGrid},
  hline{6,10,13} = {0.4pt, msGrid},
  row{13} = {font=\bfseries},
}
Question & Scheme & Marks & Notes \\
6(a) & {\((\alpha-\beta)^2+(\beta-\gamma)^2+(\gamma-\alpha)^2\)\\[4pt] \(\phantom{A}= 2\sum\alpha^2 - 2\sum\alpha\beta\)\\[4pt] \(\phantom{A}= 2\big((\sum\alpha)^2 - 2\sum\alpha\beta\big) - 2\sum\alpha\beta\)} & M1 & {Expands the squares and writes the expression in terms of \(\sum\alpha\) and \(\sum\alpha\beta\) only. Condone sign and coefficient slips.} \\
& {\(\phantom{A}= 2\big((\sum\alpha)^2 - 3\sum\alpha\beta\big)\)} & A1 & {oe, such as \(2(\sum\alpha)^2 - 6\sum\alpha\beta\).} \\
& {\(\sum\alpha = 5,\ \sum\alpha\beta = p\)\\[4pt] \(14 = 2(5^2 - 3p) \Rightarrow p = 6\)} & A1~AG & {Correct \(\sum\alpha = 5\) and \(\sum\alpha\beta = p\) substituted into a correct expression, leading without error to \(p = 6\).} \\
\nopagebreak & & [3] & \\
6(b) & {\(\alpha\beta\gamma = q\)} & B1 & \\
& {\(\frac{\sum\alpha\beta}{\alpha\beta\gamma} = \frac{6}{q} = 3\)} & M1 & {Writes the reciprocal sum as \(\sum\alpha\beta\) over \(\alpha\beta\gamma\), substitutes their product and \(p\), and sets it equal to 3. Condone a sign slip in the product.} \\
& {\(q = 2\)} & A1 & \\
\nopagebreak & & [3] & \\
6(c) & {\((\alpha+1)(\beta+1)(\gamma+1)\)\\[4pt] \(\phantom{A}= \alpha\beta\gamma + \sum\alpha\beta + \sum\alpha + 1\)\\[4pt] \(\phantom{A}= \their{2} + 6 + 5 + 1\)} & M1 & {Expands the product into \(\alpha\beta\gamma\), \(\sum\alpha\beta\), \(\sum\alpha\) and 1, condoning one slip, and substitutes their values; or any method not finding the roots, eg \(-f(-1)\).} \\
& {\(\phantom{A}= 14\)} & A1 & {cao. Scores 0 if the roots are found or no working is shown.} \\
\nopagebreak & & [2] & \\
\nopagebreak \SetCell[c=2]{r} Total for Question 6 & & 8 & \\
\end{msscheme}

\msQuestion{7}
\begin{msscheme}{
  hline{3-5} = {2-Z}{0.4pt, msGrid},
  hline{6} = {0.4pt, msGrid},
  row{6} = {font=\bfseries},
}
Question & Scheme & Marks & Notes \\
7 & {\(\alpha+\beta+\gamma=4,\ \sum\alpha\beta=-1,\ \alpha\beta\gamma=-6\)} & B1 & {All three correct, soi.} \\
& {Single fraction, then root sums:\\[4pt] \(\dfrac{\sum(\beta+\gamma)(\gamma+\alpha)}{(\alpha+\beta)(\beta+\gamma)(\gamma+\alpha)} = \ldots\)} & M1 & {Forms a single fraction and attempts numerator and denominator in \(\sum\alpha\), \(\sum\alpha\beta\), \(\alpha\beta\gamma\), or via \(\alpha+\beta=4-\gamma\) oe. Condone expansion slips.} \\
& {\(\phantom{A}=\dfrac{(\sum\alpha)^2+\sum\alpha\beta}{\sum\alpha\sum\alpha\beta-\alpha\beta\gamma}\)} & A1 & {Correct numerator and denominator, oe, or values 15 and 2 found.} \\
& {\(\phantom{A}=\dfrac{16-1}{-4+6}=\dfrac{15}{2}\)} & A1 & {\(\frac{15}{2}\) oe.} \\
\nopagebreak \SetCell[c=2]{r} Total for Question 7 & & 4 & \\
\end{msscheme}

\msQuestion{8}
\begin{msscheme}{
  hline{3-6,8-10,12-14} = {2-Z}{0.4pt, msGrid},
  hline{7,11,15} = {0.4pt, msGrid},
  row{15} = {font=\bfseries},
}
Question & Scheme & Marks & Notes \\
8(a) & {\(y = \frac{1}{x+1} \Rightarrow x = \frac{1-y}{y}\)} & B1 & {oe, such as \(x = \frac1y - 1\).} \\
& {\(\left(\frac{1-y}{y}\right)^3 + \left(\frac{1-y}{y}\right)^2\)\\[4pt] \(\phantom{A} - 2\left(\frac{1-y}{y}\right) - 1 = 0\)} & M1 & {Substitutes their \(x\) in terms of \(y\) into every term of the given cubic.} \\
& {\((1-y)^3 + y(1-y)^2 - 2y^2(1-y) - y^3 = 0\)\\[4pt] \(\phantom{A} \Rightarrow \ldots\)} & M1 & {Clears the fractions and expands to a polynomial in \(y\). Condone slips in expansion.} \\
& {\(y^3 - y^2 - 2y + 1 = 0\)} & A1 & {Or any multiple, as an equation in any variable.} \\
\nopagebreak & & [4] & \\
8(b) & {Sum of roots \(= 1\), sum of pairs \(= -2\)} & B1 & {Both values.} \\
& {\(\their{1}^2 - 2(\their{{-2}}) = \ldots\)} & M1 & {Uses (sum of roots)\(^2 - 2\)(sum of pairs) with their values.} \\
& {\(= 5\)} & A1 & {With no incorrect values seen.} \\
\nopagebreak & & [3] & \\
8(c) & {Product of roots \(= -1\)} & B1 & \\
& {\(\their{1}^3 - 3(\their{1})(\their{{-2}}) + 3(\their{{-1}}) = \ldots\)} & M1 & {Uses a correct identity for the sum of cubes, or sums the cubic over its roots, with their values. Not for the identity alone.} \\
& {\(= 4\)} & A1 & \\
\nopagebreak & & [3] & \\
\nopagebreak \SetCell[c=2]{r} Total for Question 8 & & 10 & \\
\end{msscheme}

\msQuestion{9}
\begin{msscheme}{
  hline{3-7,9} = {2-Z}{0.4pt, msGrid},
  hline{8,10} = {0.4pt, msGrid},
  row{10} = {font=\bfseries},
}
Question & Scheme & Marks & Notes \\
9(a) & {\(\frac{a}{r} + a + ar = 14\)} & M1 A1 & {\msNote{M1}{Writes the roots as a geometric progression, such as \(\frac{a}{r}, a, ar\) or \(a, ar, ar^2\), in the sum of roots. Condone a sign error in 14.}\\ \msNote{A1}{Correct sum relation, oe.}} \\
& {\(\frac{a^2}{r} + a^2 + a^2 r = 56\)\\[4pt] \(\Rightarrow a^2\left(\frac1r + 1 + r\right) = 56 \Rightarrow a = \ldots\)} & M1 & {Forms the sum of products in pairs equal to 56 in their progression and eliminates \(r\), e.g. by dividing by the sum relation. Condone slips.} \\
& {\(a = 4\)} & A1 & {Middle root 4, e.g. \(ar = 4\) for roots \(a, ar, ar^2\).} \\
& {\(\their{4}\left(\frac1r + 1 + r\right) = 14\)\\[4pt] \(\Rightarrow 2r^2 - 5r + 2 = 0 \Rightarrow r = \ldots\ [= \tfrac12, 2]\)} & dM1 & {Substitutes their middle root into a relation to form a three-term quadratic in \(r\) and solves it by any means, leading to \(r = \ldots\)} \\
& {Roots \(2,\ 4,\ 8\)} & A1 & {All three roots stated; values of \(r\) alone do not score.} \\
\nopagebreak & & [6] & \\
9(b) & {\(c = -(2 \times 4 \times 8) = -64\)} & B1ft & {ft minus the product of their three roots.} \\
\nopagebreak & & [1] & \\
\nopagebreak \SetCell[c=2]{r} Total for Question 9 & & 7 & \\
\end{msscheme}

\msQuestion{10}
\begin{msscheme}{
  hline{3-5,7-9,11-13} = {2-Z}{0.4pt, msGrid},
  hline{6,10,14} = {0.4pt, msGrid},
  row{14} = {font=\bfseries},
}
Question & Scheme & Marks & Notes \\
10(a) & {\(\sum\alpha\beta\gamma = 4,\quad \alpha\beta\gamma\delta = 1\)} & B1 & \\
& {\(\frac1\alpha+\frac1\beta+\frac1\gamma+\frac1\delta = \frac{\sum\alpha\beta\gamma}{\alpha\beta\gamma\delta} = \ldots\)} & M1 & {Writes the reciprocal sum as their triple-product sum over their product of roots, or uses the sum of roots of the reciprocal equation.} \\
& {\(= 4\)} & A1 & \\
\nopagebreak & & [3] & \\
10(b) & {\(\sum\frac1{\alpha\beta} = \frac{\sum\alpha\beta}{\alpha\beta\gamma\delta} = 3\)} & B1 & \\
& {\(\sum\frac1{\alpha^2} = \left(\sum\frac1\alpha\right)^2 - 2\sum\frac1{\alpha\beta}\)\\[4pt] \(\phantom{\sum\frac1{\alpha^2}} = \their{4}^2 - 2(\their{3}) = \ldots\)} & M1 & {Uses the square of their reciprocal sum minus twice their pair sum of reciprocals.} \\
& {\(= 10\)} & A1 & \\
\nopagebreak & & [3] & \\
10(c) & {\(\frac1\alpha+\frac1\beta+\frac1\gamma = 4-\frac1\delta\)\\[4pt] \(\frac1\beta+\frac1\gamma+\frac1\delta = 4-\frac1\alpha\)\\[4pt] \(\frac1\gamma+\frac1\delta+\frac1\alpha = 4-\frac1\beta\)\\[4pt] \(\frac1\delta+\frac1\alpha+\frac1\beta = 4-\frac1\gamma\)} & M1 A1 & {\msNote{M1}{Writes at least one bracket as the sum of all four reciprocals minus the omitted reciprocal, or expands one bracket squared into its three squares and twice its pair products.}\\ \msNote{A1}{All four correct, or that bracket's expansion correct.}} \\
& {\(\sum\left(4-\frac1\alpha\right)^2 = \sum\left(16-\frac8\alpha+\frac1{\alpha^2}\right)\)\\[4pt] \(\phantom{\sum} = 64 - 8\sum\frac1\alpha + \sum\frac1{\alpha^2}\)} & M1 A1 & {\msNote{M1}{Expands the squared brackets and sums, collecting in terms of \(\sum\frac1\alpha\) and \(\sum\frac1{\alpha^2}\), or of \(\sum\frac1{\alpha^2}\) and \(\sum\frac1{\alpha\beta}\). Condone coefficient slips.}\\ \msNote{A1}{Correct collected form, oe such as \(3\sum\frac1{\alpha^2}+4\sum\frac1{\alpha\beta}\).}} \\
& {\(= 64 - 32 + 10 = 42\)} & A1 & \\
\nopagebreak & & [5] & \\
\nopagebreak \SetCell[c=2]{r} Total for Question 10 & & 11 & \\
\end{msscheme}

\msQuestion{11}
\begin{msscheme}{
  hline{3-6,8-10,12-13} = {2-Z}{0.4pt, msGrid},
  hline{7,11,14} = {0.4pt, msGrid},
  row{14} = {font=\bfseries},
}
Question & Scheme & Marks & Notes \\
11(a) & {\(t = \frac{1}{x-1} \Rightarrow x = 1 + \frac{1}{t}\)} & B1 & {Or \(x = \frac{1+t}{t}\) oe; any letter.} \\
& {\(\left(1+\frac{1}{t}\right)^3 - 3\left(1+\frac{1}{t}\right) + 1 = 0\)} & M1 & {Substitutes their \(x\) in terms of \(t\) into the cubic.} \\
& {\(\frac{1}{t^3} + \frac{3}{t^2} - 1 = 0 \Rightarrow 1 + 3t - t^3 = 0\)} & M1 & {Expands the cube and multiplies through by \(t^3\) to reach a polynomial equation. Condone slips in the expansion.} \\
& {\(t^3 - 3t - 1 = 0\)} & A1 & {Or any multiple, in any variable; must be an equation.} \\
\nopagebreak & & [4] & \\
11(b) & {\(\sum r = 0,\ \sum r_1r_2 = -3\)} & B1 & {Both values, from the cubic in (a).} \\
& {\(\sum r^2 = \their{0}^2 - 2(\their{-3})\)} & M1 & {Uses \((\sum r)^2 - 2\sum r_1r_2\) with their values from a cubic in the reciprocal variable.} \\
& {\(= 6\)} & A1 & \\
\nopagebreak & & [3] & \\
11(c) & {\(r^3 = 3r + 1 \Rightarrow r^5 = 3r^3 + r^2\)\\[4pt] \(\phantom{r^3 = 3r + 1 \Rightarrow r^5} = r^2 + 9r + 3\)} & M1 A1 & {\msNote{M1}{Uses \(r^3 = 3r + 1\) to express \(r^5\) in powers below 3 and sums over the roots, or uses the recurrence \(\sum r^{n+3} = 3\sum r^{n+1} + \sum r^n\). Condone slips.}\\ \msNote{A1}{oe, such as \(\sum r^5 = 3\sum r^3 + \sum r^2\) with \(\sum r^3 = 3\).}} \\
& {\(\sum r^5 = \their{6} + 9(0) + 3 \times 3 = 15\)} & A1 & \\
\nopagebreak & & [3] & \\
\nopagebreak \SetCell[c=2]{r} Total for Question 11 & & 10 & \\
\end{msscheme}

\msQuestion{12}
\begin{msscheme}{
  hline{3-5,7-10} = {2-Z}{0.4pt, msGrid},
  hline{6,11} = {0.4pt, msGrid},
  row{11} = {font=\bfseries},
}
Question & Scheme & Marks & Notes \\
12(a) & {\(x = \frac1y\)} & B1 & {Or \(y = \frac1x\); any letter.} \\
& {\(\frac{3}{y^4} - \frac{1}{y^3} - \frac{10}{y^2} + \frac{8}{y} + 4 = 0\)\\[4pt] \(\Rightarrow 3 - y - 10y^2 + 8y^3 + 4y^4 = 0\)} & M1 & {Substitutes a reciprocal into every term and multiplies through by a power of \(y\) to clear fractions. Condone sign and coefficient slips.} \\
& {\(4x^4 + 8x^3 - 10x^2 - x + 3 = 0\)} & A1 & {Or any non-zero integer multiple; any variable.} \\
\nopagebreak & & [3] & \\
12(b) & {\(x = 1 + \frac2y\)} & B1 & {Or \(y = \frac{2}{x-1}\), or \(x = \frac{y+2}{y}\); any letter.} \\
& {\(3\left(1+\frac2y\right)^4 - \left(1+\frac2y\right)^3 - 10\left(1+\frac2y\right)^2\)\\[4pt] \(\phantom{3} + 8\left(1+\frac2y\right) + 4 = 0\)} & M1 & {Substitutes their \(x\) of the form \(1 + \frac{k}{y}\), oe, into every term of the quartic.} \\
& {\(3(y+2)^4 - y(y+2)^3 - 10y^2(y+2)^2\)\\[4pt] \(\phantom{3} + 8y^3(y+2) + 4y^4 = 0 \Rightarrow \ldots\)} & M1 & {Expands every power, before or after multiplying through by \(y^4\), and collects to a polynomial in \(y\). Condone arithmetic slips.} \\
& {\(2x^4 - 3x^3 + 10x^2 + 44x + 24 = 0\)} & A1 & {Or any non-zero integer multiple, such as \(4x^4 - 6x^3 + 20x^2 + 88x + 48 = 0\); any variable.} \\
\nopagebreak & & [4] & \\
\nopagebreak \SetCell[c=2]{r} Total for Question 12 & & 7 & \\
\end{msscheme}

\msQuestion{13}
\begin{msscheme}{
  hline{3-8,10-13} = {2-Z}{0.4pt, msGrid},
  hline{9,14} = {0.4pt, msGrid},
  row{14} = {font=\bfseries},
}
Question & Scheme & Marks & Notes \\
13(a) & {\(\alpha+\beta+\gamma+\delta = -b \Rightarrow b = 2\)} & B1 & \\
& {\(\sum\frac{1}{\alpha} = \frac{\sum\alpha\beta\gamma}{\alpha\beta\gamma\delta} = \frac{-d}{3}\)\\[4pt] \(\phantom{A}\Rightarrow \frac{-d}{3} = 2 \Rightarrow d = \ldots\)} & M1 & {Writes the reciprocal sum as the sum of triple products over the product of roots, in terms of \(d\) and 3, and solves for \(d\). Condone a sign slip.} \\
& {\(d = -6\)} & A1 & \\
& {\(\sum\alpha^2 + 2\sum\alpha + 4 = 6\)\\[4pt] \((\sum\alpha)^2 = \sum\alpha^2 + 2\sum\alpha\beta\)} & M1 & {Expands the squared condition into \(\sum\alpha^2\), \(\sum\alpha\) and a constant, and relates \(\sum\alpha^2\) to \((\sum\alpha)^2\) and \(c\). Condone slips in coefficients or the constant.} \\
& {\(\sum\alpha^2 = 6\)} & A1 & {soi, eg within the equation for \(c\).} \\
& {\((-2)^2 = 6 + 2c \Rightarrow c = -1\)} & A1 & \\
\nopagebreak & & [6] & \\
13(b) & {\(\sum\alpha^3 + 3\sum\alpha^2 + 3\sum\alpha + 4 = 20\)} & M1 & {Expands the cubed condition into \(\sum\alpha^3\), \(\sum\alpha^2\), \(\sum\alpha\) and a constant. Condone slips in the coefficients or the constant.} \\
& {\(\sum\alpha^3 + 18 - 6 + 4 = 20\)\\[4pt] \(\phantom{A}\Rightarrow \sum\alpha^3 = 4\)} & A1 & \\
& {\(\sum\alpha^4 = -2\sum\alpha^3 + \sum\alpha^2 + 6\sum\alpha - 12\)\\[4pt] \(\phantom{\sum\alpha^4} = -2(\their{4}) + 6 + 6(-2) - 12\)} & M1 & {Expresses \(\sum\alpha^4\) in their lower power sums, by summing their quartic over the roots or otherwise, and substitutes. Condone a constant of \(-3\) for \(-12\).} \\
& {\(\sum\alpha^4 = -26\)} & A1 & \\
\nopagebreak & & [4] & \\
\nopagebreak \SetCell[c=2]{r} Total for Question 13 & & 10 & \\
\end{msscheme}

\msQuestion{14}
\begin{msscheme}{
  hline{3-7,9-11,13-14} = {2-Z}{0.4pt, msGrid},
  hline{8,12,15} = {0.4pt, msGrid},
  row{15} = {font=\bfseries},
}
Question & Scheme & Marks & Notes \\
14(a) & {\(y = x^2\)} & B1 & {States or uses \(y = x^2\), or \(x = \sqrt{y}\).} \\
& {\((x^2-1)^2 = -2x\)\\[4pt] \(\Rightarrow (y-1)^2 = -2\sqrt{y}\)} & M1 & {Substitutes \(x = \sqrt{y}\) with the \(\sqrt{y}\) term separated from the integer powers of \(y\). Condone sign slips.} \\
& {\((y-1)^4 = 4y\)\\[4pt] \(\Rightarrow y^4-4y^3+6y^2-4y+1 = 4y\)} & M1 & {Squares both sides to eliminate \(\sqrt{y}\) and expands to a polynomial in \(y\). Condone sign and coefficient slips.} \\
& {\(y^4-4y^3+6y^2-8y+1=0\)} & A1 & {Any non-zero multiple, in any variable; must be an equation.} \\
& {\(\alpha^2+\beta^2+\gamma^2+\delta^2 = 4\)} & B1ft & {ft minus the ratio of the cubic to quartic coefficients of their quartic; or 4 from \((\sum\alpha)^2 - 2\sum\alpha\beta\).} \\
\nopagebreak & & [5] & \\
14(b) & {\(\sum\alpha^2 = 4,\ \sum\alpha^2\beta^2 = 6,\ \sum\alpha^2\beta^2\gamma^2 = 8\)} & B1 & {All three values, from the quartic in \(y\) or from the original roots.} \\
& {\(\sum\alpha^6 = 4^3 - 3(4)(6) + 3(8)\)} & M1 & {Uses a correct identity for the sum of cubes in terms of the root sums, with their values; or any valid power-sum method.} \\
& {\(\phantom{\sum\alpha^6} = 16\)} & A1 & \\
\nopagebreak & & [3] & \\
14(c) & {\(\sum\alpha^4 = 4^2 - 2(6)\ [=4]\)\\[4pt] \(\sum\alpha^8 = 4\sum\alpha^6 - 6\sum\alpha^4 + 8\sum\alpha^2 - 4\)\\[4pt] \(\phantom{\sum\alpha^8} = 4(\their{16}) - 6(4) + 8(4) - 4\)} & M1 & {Expresses \(\sum\alpha^8\) in lower power sums via the quartic in \(y\) or another valid identity, finding \(\sum\alpha^4\) and substituting their values. Condone slips.} \\
& {\(\phantom{\sum\alpha^8} = 68\)} & A1 & \\
\nopagebreak & & [2] & \\
\nopagebreak \SetCell[c=2]{r} Total for Question 14 & & 10 & \\
\end{msscheme}

\msQuestion{15}
\begin{msscheme}{
  hline{3-5,7-9} = {2-Z}{0.4pt, msGrid},
  hline{6,10} = {0.4pt, msGrid},
  row{10} = {font=\bfseries},
}
Question & Scheme & Marks & Notes \\
15(a) & {\(\left(z^2+\frac{1}{z^2}\right)-2\left(z+\frac{1}{z}\right)-1=0\)\\[4pt] \(\Rightarrow (w^2-2)-2w-1=0\)} & M1 & {Divides by \(z^2\) and uses \(z^2+\frac{1}{z^2}=w^2\pm2\) to form an equation in \(w\) only. Condone sign and coefficient slips.} \\
& {\(w^2-2w-3=0\)} & A1 & \\
& {Roots in reciprocal pairs \(\Rightarrow w = 3, 3, -1, -1\)\\[4pt] \((w^2-2w-3)^2=0\)\\[4pt] \(\Rightarrow w^4-4w^3-2w^2+12w+9=0\)} & A1~AG & {States that roots occur in reciprocal pairs, or that each value of \(w\) arises twice, and reaches the printed quartic without error. A0 for squaring with no reason given.} \\
\nopagebreak & & [3] & \\
15(b) & {\(z+\frac{1}{z}=3 \Rightarrow z^2-3z+1=0\)\\[4pt] \(z+\frac{1}{z}=-1 \Rightarrow z^2+z+1=0\)\\[4pt] \(\Rightarrow z=\ldots\)} & M1 & {Forms \(z^2-wz+1=0\) for both values of \(w\) from the quartic and solves each by any means. Condone sign slips.} \\
& {\(z=\frac{3\pm\sqrt5}{2}\)} & A1 & {Both, exact, oe.} \\
& {\(z=\frac{-1\pm i\sqrt3}{2}\)} & A1 & {Both, exact, oe such as \(e^{\pm 2\pi i/3}\).} \\
\nopagebreak & & [3] & \\
\nopagebreak \SetCell[c=2]{r} Total for Question 15 & & 6 & \\
\end{msscheme}

\end{document}
